\documentclass{article}
\usepackage[utf8]{inputenc}
\usepackage{amsmath,amssymb}
\usepackage[most]{tcolorbox}
\usepackage{geometry}
\geometry{margin=2.5cm}

\newcommand{\step}[1]{\par\noindent\hangindent=3.5em\hangafter=1%
  \makebox[3.5em][l]{\textbf{#1.}}}
\newcommand{\steptag}[1]{\hfill{\small\textsf{[#1]}}}

\newtcolorbox{proofbox}[1]{
  colback=gray!5, colframe=gray!60,
  fonttitle=\bfseries, title={#1},
  breakable, enhanced,
  left=4pt, right=4pt, top=4pt, bottom=4pt,
  before skip=10pt, after skip=10pt
}

\begin{document}

\begin{proofbox}{Bound extra fixed class with one ledger selection: there ...}
\step{1} Bound extra fixed class with one ledger selection: there exist one def-stage1-polar-witness-data tuple W=(C\_rect,C\_ch,C\_pol,C\_grp,C\_path,C\_der,e\_rect,kappa\_ch,kappa\_pol,kappa\_der,delta\_*,epsilon\_*\textasciicircum{}r,e\_S1,r\_iso) such that C\_rect,C\_ch,C\_pol,C\_grp,C\_path,C\_der>=1, 0<e\_rect<=1/C\_rect, 0<kappa\_ch,kappa\_pol,kappa\_der<=1/2, delta\_*=min\{1/4,kappa\_ch/(4*C\_ch),kappa\_pol/(4*C\_pol)\}, epsilon\_*\textasciicircum{}r=min\{1/4,kappa\_ch/(4*C\_ch),kappa\_pol/(4*C\_pol),kappa\_der/(8*C\_der),1/C\_grp,delta\_*/(12*C\_path*C\_grp)\}, e\_S1=min\{e\_rect,epsilon\_*\textasciicircum{}r/C\_rect\}, r\_iso=min\{delta\_*/4,kappa\_der/(8*C\_der)\}, and there are universal e\_quot\textasciicircum{}r>0, epsilon\_B\textasciicircum{}r>0, C\_fix<infinity, 0<e\_fix\textasciicircum{}r<=epsilon\_B\textasciicircum{}r, and r\_bidx>0 with epsilon\_B\textasciicircum{}r=min\{epsilon\_*\textasciicircum{}r,e\_quot\textasciicircum{}r\} and C\_fix>=C\_grp such that, for every finite-dimensional exact-unit epsilon\_r-C*-algebra with 0<=epsilon\_r<=e\_fix\textasciicircum{}r and 1<dim\_C calX<infinity, there exist a family g=(g\_V:B\_\{2delta\_*\}\textasciicircum{}\{icalH\}(0)->B\_\{2delta\_*\}\textasciicircum{}\{calH\}(0))\_\{V in calUbar\_delta\_*\}, maps u\_delta,h\_delta,mu,sigma,H, a space breve-calU, maps breve-mu:breve-calU x breve-calU->breve-calU and breve-sigma:breve-calU->breve-calU, a breve-sigma-fixed class breve-U, U\_0 in calU\_e, c,a in U(1), and U in calU\_e such that all of the following hold: (A\_2) for every V in calUbar\_delta\_* and A\textasciicircum{}par in B\_\{2delta\_*\}\textasciicircum{}\{icalH\}(0), g\_V(A\textasciicircum{}par) is the unique element of B\_\{2delta\_*\}\textasciicircum{}\{calH\}(0) with f\_V(A\textasciicircum{}par+g\_V(A\textasciicircum{}par))=0, where f\_V(A)=(1/2)*(((J+A\textasciicircum{}dagger) bold-dot V\textasciicircum{}dagger) bold-dot (V bold-dot (J+A))-J), chi\_V(A\textasciicircum{}par)=V bold-dot (J+A\textasciicircum{}par+g\_V(A\textasciicircum{}par)) lies in calU, ||Dg\_V(A\textasciicircum{}par)||<=C\_ch*(epsilon\_r+delta\_*), ||D\_\{A\textasciicircum{}perp\}f\_V(A\textasciicircum{}par+g\_V(A\textasciicircum{}par))-I\_calH||<=C\_ch*(epsilon\_r+delta\_*)<1, and these C\textasciicircum{}1 charts cover calU; (A\_4) set Pi\_delta\_*:calU x B\_\{delta\_*\}\textasciicircum{}\{calH\}(J)->calX by Pi\_delta\_*(V,K):=V bold-dot K and set S\_delta\_*:=Pi\_delta\_*(calU x B\_\{delta\_*\}\textasciicircum{}\{calH\}(J)); Pi\_delta\_* is a C\textasciicircum{}1 diffeomorphism onto S\_delta\_* with unique inverse maps u\_delta:S\_delta\_*->calU and h\_delta:S\_delta\_*->B\_\{delta\_*\}\textasciicircum{}\{calH\}(J) satisfying X=u\_delta(X) bold-dot h\_delta(X) for every X in S\_delta\_* and u\_delta(V)=V and h\_delta(V)=J for every V in calU; (A\_5) the displayed maps mu:calU x calU->calU and sigma:calU->calU are defined by mu(V\_1,V\_2):=u\_delta(V\_1 bold-dot V\_2) and sigma(V):=u\_delta(V\textasciicircum{}dagger), are global C\textasciicircum{}1 and, for every V,V\_1,V\_2,V\_3 in calU, mu(J,V)=mu(V,J)=V, sigma(J)=J, ||mu(V\_1,V\_2)-V\_1 bold-dot V\_2||<=C\_grp*epsilon\_r, ||sigma(V)-V\textasciicircum{}dagger||<=C\_grp*epsilon\_r, ||mu(mu(V\_1,V\_2),V\_3)-mu(V\_1,mu(V\_2,V\_3))||<=C\_grp*epsilon\_r, ||mu(sigma(V),V)-J||<=C\_grp*epsilon\_r, and ||mu(V,sigma(V))-J||<=C\_grp*epsilon\_r; (A\_6) for every V\_0,V\_1 in calU and q in [0,1] satisfying ||V\_1-V\_0||<=q, C\_path*q<=1/4, and C\_path*(q+epsilon\_r*q+q\textasciicircum{}2)<delta\_*-C\_pol*(epsilon\_r*delta\_*+delta\_*\textasciicircum{}2), the path H(-,V\_0,V\_1):[0,1]->calU given by H(t,V\_0,V\_1):=u\_delta((1-t)*V\_0+t*V\_1) is defined, is jointly continuous in its displayed variables, and joins V\_0 to V\_1, and satisfies H(t,c\_0*V\_0,c\_0*V\_1)=c\_0*H(t,V\_0,V\_1) for every c\_0 in U(1) and t in [0,1]; (A\_7) for every s in \{+1,-1\}, set chi\_s:B\_\{r\_iso\}\textasciicircum{}\{icalH\}(0)->calU by chi\_s(A):=sJ bold-dot (J+A+g\_\{sJ\}(A)), let phi\_\{sJ\}\textasciicircum{}par:chi\_s(B\_\{r\_iso\}\textasciicircum{}\{icalH\}(0))->B\_\{r\_iso\}\textasciicircum{}\{icalH\}(0) be its inverse, and set F\_s:B\_\{r\_iso\}\textasciicircum{}\{icalH\}(0)->icalH by F\_s(A):=phi\_\{sJ\}\textasciicircum{}par(sigma(chi\_s(A))); then sigma maps chi\_s(B\_\{r\_iso\}\textasciicircum{}\{icalH\}(0)) into itself and ||D(F\_s-id)(A)+2I||<=C\_der*(epsilon\_r+r\_iso) for every A in B\_\{r\_iso\}\textasciicircum{}\{icalH\}(0); (R) set q\_*:=C\_grp*epsilon\_r, set r\_-:=delta\_*-C\_pol*(epsilon\_r*delta\_*+delta\_*\textasciicircum{}2), and set eta\_*:=C\_path*(q\_*+epsilon\_r*q\_*+q\_*\textasciicircum{}2); the guards C\_ch*(epsilon\_r+delta\_*)<=kappa\_ch, C\_pol*(epsilon\_r+delta\_*)<=kappa\_pol, q\_*<r\_-, C\_path*q\_*<=1/4, eta\_*<r\_-, C\_der*(epsilon\_r+r\_iso)<=kappa\_der/4<1, and (1+epsilon\_r)*(1+C\_ch*(epsilon\_r+delta\_*))*r\_iso+q\_*<2delta\_* hold; moreover r\_bidx=r\_iso and breve-calU=calU\_e/U(1), set breve-e:=[J], breve-mu([V\_1],[V\_2])=[mu(V\_1,V\_2)], breve-sigma([V])=[sigma(V)], (breve-calU,breve-mu,breve-e) is a connected H-space, breve-sigma is a smooth left inversion, sigma(c\_0*V)=conj(c\_0)*sigma(V) and ||sigma(V)-V\textasciicircum{}dagger||<=C\_grp*epsilon\_r for every c\_0 in U(1) and V in calU, breve-calU is a connected compact orientable smooth manifold without boundary of real dimension (dim\_C calX)-1 and is homeomorphic to a finite simplicial complex, breve-e is an isolated fixed point of breve-sigma with local index +1, J and -J are the only sigma-fixed points in their respective ambient r\_bidx-balls, breve-U!=breve-e, [U\_0]=breve-U, sigma(U\_0)=c*U\_0, a\textasciicircum{}2=c, U=a*U\_0, [U]=breve-U, sigma(U)=U, ||U-U\textasciicircum{}dagger||<=C\_fix*epsilon\_r, ||U-J||>=r\_bidx, and ||U+J||>=r\_bidx. \steptag{validated} \\[6pt]
\step{1.1} One-use ledger selection. By lem-stage1-polar-constant-ledger there exists one universal def-stage1-polar-witness-data tuple W with the six coefficients at least 1, the stated positive margin bounds, and exactly the displayed formulas for delta\_*, epsilon\_*\textasciicircum{}r, e\_S1, and r\_iso; for this same W its clauses (A\_2),(A\_4),(A\_5),(A\_6),(A\_7) are available for every exact-unit epsilon\_r-C*-algebra whenever their own displayed scalar guards hold. This is the sole existential elimination of lem-stage1-polar-constant-ledger, and its rectified-input clause (R) is not used for the arbitrary exact-unit input below. \steptag{validated} \\[4pt]
\step{1.2} Conditional bound-package application. For every W satisfying the hypotheses and formulas in node 1.1, lem-stage1-bound-quotient-index-data gives universal e\_quot\textasciicircum{}r>0 and epsilon\_B\textasciicircum{}r=min\{epsilon\_*\textasciicircum{}r,e\_quot\textasciicircum{}r\}>0 such that, whenever a finite-dimensional exact-unit epsilon\_r-C*-algebra with epsilon\_r<=epsilon\_B\textasciicircum{}r and 1<dim\_C calX<infinity is supplied with the literal (A\_2),(A\_4),(A\_5),(A\_6),(A\_7) data and the listed scalar guards, it returns those same g,u\_delta,h\_delta,mu,sigma,H together with r\_bidx=r\_iso, breve-calU=calU\_e/U(1), breve-mu, breve-sigma, the connected compact orientable smooth boundaryless (dim\_C calX-1)-manifold and finite-simplicial-complex package, the H-space and smooth-left-inversion assertions, scalar covariance and near-adjoint estimate, local index +1 at the isolated breve-e=[J], both ambient sigma-isolation assertions, and the class-first phase-lift witnesses for every breve-sigma-fixed class. \steptag{validated} \\[4pt]
\step{1.3} Scalar guards for arbitrary exact-unit inputs. Fix W as in node 1.1 and 0<=epsilon\_r<=epsilon\_B\textasciicircum{}r<=epsilon\_*\textasciicircum{}r. Put q\_*=C\_grp*epsilon\_r, r\_-=delta\_*-C\_pol*(epsilon\_r*delta\_*+delta\_*\textasciicircum{}2), and eta\_*=C\_path*(q\_*+epsilon\_r*q\_*+q\_*\textasciicircum{}2). The minimum formulas imply C\_ch*epsilon\_r<=kappa\_ch/4 and C\_ch*delta\_*<=kappa\_ch/4, and analogously for C\_pol, so C\_ch*(epsilon\_r+delta\_*)<=kappa\_ch/2<=kappa\_ch and C\_pol*(epsilon\_r+delta\_*)<=kappa\_pol/2<=kappa\_pol. They also give epsilon\_r<=1/4, C\_pol*(epsilon\_r+delta\_*)<=1/4, hence r\_- >=3*delta\_*/4; q\_*<=delta\_*/(12*C\_path)<=delta\_*/12<r\_-; C\_path*q\_*<=delta\_*/12<1/4; and eta\_*=C\_path*q\_*(1+epsilon\_r+q\_*)<=61*delta\_*/576<3*delta\_*/4<=r\_-. Further C\_der*epsilon\_r<=kappa\_der/8 and C\_der*r\_iso<=kappa\_der/8, so C\_der*(epsilon\_r+r\_iso)<=kappa\_der/4<1. Finally r\_iso<=delta\_*/4, 1+epsilon\_r<=5/4, 1+C\_ch*(epsilon\_r+delta\_*)<=5/4, and q\_*<=delta\_*/12, whence (1+epsilon\_r)*(1+C\_ch*(epsilon\_r+delta\_*))*r\_iso+q\_*<=25*delta\_*/64+delta\_*/12<2*delta\_*. Thus every graph, polar, group, path, derivative, and chart-retention guard required by the conditional ledger clauses and lem-stage1-bound-quotient-index-data holds without using ledger clause (R). \steptag{validated} \\[4pt]
\step{1.4} Common-object instantiation. Choose e\_fix\textasciicircum{}r:=epsilon\_B\textasciicircum{}r and C\_fix:=C\_grp, so 0<e\_fix\textasciicircum{}r<=epsilon\_B\textasciicircum{}r, C\_fix<infinity, and C\_fix>=C\_grp. For an arbitrary finite-dimensional exact-unit epsilon\_r-C*-algebra with 0<=epsilon\_r<=e\_fix\textasciicircum{}r and 1<dim\_C calX<infinity, node 1.3 permits the conditional clauses (A\_2),(A\_4),(A\_5),(A\_6),(A\_7) from the single W selected by lem-stage1-polar-constant-ledger to be instantiated at delta\_*. Feeding exactly those functions, inverses, operations, paths, and charts to lem-stage1-bound-quotient-index-data yields all analytic clauses (A\_2),(A\_4)-(A\_7), all displayed guards (R), and the complete same-map quotient package listed in node 1.2, with r\_bidx=r\_iso. \steptag{validated} \\[4pt]
\step{1.5} Singleton reduction. For the package obtained in node 1.4, if breve-e were the only breve-sigma-fixed class, then Fix(breve-sigma)=\{breve-e\} would be a finite fixed set, and breve-e has local fixed-point index +1 by lem-stage1-bound-quotient-index-data. \steptag{validated} \\[4pt]
\step{1.6} Maximal-simplex antecedent. Under the singleton hypothesis of node 1.5, the homeomorphism of breve-calU with a finite simplicial complex makes it a finite polyhedron, and lem-finite-polyhedron-maximal-simplex-placement places its sole fixed point breve-e in a maximal simplex. Hence the maximal-simplex hypothesis of lem-topology-lefschetz-hopf is satisfied. \steptag{validated} \\[4pt]
\step{1.7} Lefschetz value under the singleton hypothesis. Apply lem-topology-lefschetz-hopf to the finite-polyhedron model of breve-sigma using node 1.6. Since the only fixed point is breve-e and its local index is +1, the Lefschetz number is L(breve-sigma)=1. The local-index value is also consistent with lem-topology-local-index-sign and the determinant data supplied by lem-stage1-bound-quotient-index-data. \steptag{validated} \\[6pt]
\step{1.7.1} Finite-polyhedron presentation without conjugating the map. Let h:breve-calU->|K| be the homeomorphism to the realization of a finite simplicial complex supplied by lem-stage1-bound-quotient-index-data, and transport the finite-polyhedron structure of |K| along h to the same underlying topological space breve-calU. Denote this finite-polyhedron presentation by X, and let f:X->X be the same underlying function breve-sigma (not h o breve-sigma o h\textasciicircum{}\{-1\}). Under the singleton hypothesis Fix(f)=\{breve-e\}; node 1.6 (using lem-finite-polyhedron-maximal-simplex-placement) places breve-e in a maximal simplex of this presentation. Because X and breve-calU have the same underlying space and f and breve-sigma are literally the same self-map, the local index occurring in lem-topology-lefschetz-hopf is exactly ind(f,breve-e)=ind(breve-sigma,breve-e)=+1 as supplied by lem-stage1-bound-quotient-index-data, and f\textasciicircum{}\{*k\}=breve-sigma\textasciicircum{}\{*k\} on the same cohomology groups, hence L(f)=L(breve-sigma). No invariance of index under homeomorphic conjugacy is used. \steptag{validated} \\[4pt]
\step{1.7.2} Direct Lefschetz-Hopf application, without conjugating the map. Under the singleton hypothesis, validated node 1.5 gives Fix(breve-sigma)=\{breve-e\} and validated node 1.6 gives that breve-calU, with the finite-polyhedron presentation supplied by the bound quotient package, is a finite polyhedron whose unique fixed point breve-e lies in a maximal simplex. The same package instantiated in validated node 1.4 gives ind(breve-sigma,breve-e)=+1. Therefore lem-topology-lefschetz-hopf applies directly to breve-sigma:breve-calU->breve-calU and yields L(breve-sigma)=sum\_\{x in Fix(breve-sigma)\} ind(breve-sigma,x)=ind(breve-sigma,breve-e)=1. No conjugated map f and no homeomorphic-conjugacy invariance of local fixed-point index are used. \steptag{validated} \\[4pt]
\step{1.8} Trace calculation on the same quotient. The finite simplicial model gives breve-calU finite CW type and finite-dimensional total real cohomology; the package makes it connected and an H-space and makes breve-sigma a left inversion. Therefore lem-stage1-left-inversion-trace applies and gives Tr(breve-sigma\textasciicircum{}\{*k\}:H\textasciicircum{}k(breve-calU;reals)->H\textasciicircum{}k(breve-calU;reals))=(-1)\textasciicircum{}k b\_k for every k>=0, where b\_k=dim\_reals H\textasciicircum{}k(breve-calU;reals). \steptag{validated} \\[4pt]
\step{1.9} Lefschetz number as total Betti number. By the definition in def-lefschetz-fixed-point-data and node 1.8, L(breve-sigma)=sum\_k (-1)\textasciicircum{}k Tr(breve-sigma\textasciicircum{}\{*k\})=sum\_k (-1)\textasciicircum{}k(-1)\textasciicircum{}k b\_k=sum\_k b\_k; the sum is finite because breve-calU has finite CW type. \steptag{validated} \\[4pt]
\step{1.10} Existence of an extra fixed class. Let d=dim\_R breve-calU=dim\_C calX-1>0. Connectedness gives b\_0=1, while lem-topology-orientable-top-cohomology applied to the connected compact orientable boundaryless d-manifold gives H\textasciicircum{}d(breve-calU;reals)!=0 and hence b\_d>=1. Since d>0, node 1.9 yields L(breve-sigma)=sum\_k b\_k>=2, contradicting node 1.7 under the singleton hypothesis. Thus breve-sigma has a fixed class breve-U!=breve-e. \steptag{validated} \\[4pt]
\step{1.11} Fixed lift and near-adjoint bound. Apply the class-first phase clause of lem-stage1-bound-quotient-index-data to the extra fixed class from node 1.10: bind U\_0 in calU\_e and c,a in U(1) with [U\_0]=breve-U, sigma(U\_0)=c*U\_0 and a\textasciicircum{}2=c, and set U=a*U\_0. The exported covariance sigma(zV)=conj(z)*sigma(V) gives sigma(U)=conj(a)cU\_0=aU\_0=U because c=a\textasciicircum{}2 and |a|=1; also [U]=[U\_0]=breve-U and U lies in calU\_e. The literal package estimate then gives ||U-U\textasciicircum{}dagger||=||sigma(U)-U\textasciicircum{}dagger||<=C\_grp*epsilon\_r<=C\_fix*epsilon\_r. \steptag{validated} \\[4pt]
\step{1.12} Distance bounds and final assembly. The lift U in node 1.11 is sigma-fixed and [U]=breve-U!=breve-e=[J]. If ||U-J||<r\_bidx, the J-ball isolation from lem-stage1-bound-quotient-index-data forces U=J, contradicting its class; hence ||U-J||>=r\_bidx. If ||U+J||<r\_bidx, the -J-ball isolation forces U=-J, whose scalar class equals [J], again a contradiction; hence ||U+J||>=r\_bidx. Combining nodes 1.1-1.4 and 1.10-1.11 with e\_fix\textasciicircum{}r=epsilon\_B\textasciicircum{}r, C\_fix=C\_grp, and r\_bidx=r\_iso supplies every existential witness, formula, analytic clause, guard, topological assertion, fixed lift, estimate, and distance conclusion in root node 1. \steptag{validated} \\[4pt]
\end{proofbox}

\end{document}
